% 第 19 章　电荷和电场
\chapter{电荷和电场}

\step{反常识开场}

\begin{mainline}
冬天的夜里，关了灯，你在黑暗中脱毛衣。就在毛衣离开头发的一瞬间，你听见一连串细小的“噼啪”声，眼角甚至能瞥见几点蓝白色的小火花。第二天早上梳头，头发像活过来一样，一根根立起来，追着梳子跑。开门伸手去碰金属门把手，“啪”，指尖被狠狠电了一下。

请先注意一件被我们忽略掉的怪事：这一切发生的时候，你身边没有电池，没有插座，没有任何“电器”。一件毛衣、一把塑料梳子，就是我们印象中与“电”最不相干的东西。可它们不仅能放电，还能放出光来。

再把它放大一百万倍看：夏天的雷雨夜，一道闪电劈下来，瞬间照亮半边天空。闪电是什么？我们会在本章末尾给出一个有底气的回答，这里先预告一句——闪电和你脱毛衣时的小火花，是同一种物理过程，只是规模不同。你的卧室里每晚都在上演微型雷暴，而你一直没有追问过它。

这一现象已经挑衅人类两千多年了。古希腊人发现，用毛皮摩擦琥珀，琥珀就能吸起羽毛和碎草屑。中文的“电”字本义就是闪电。可是直到大约两百年前，才有人把这些零零碎碎的怪事拼成一幅完整的图。这一章我们就重走这条路：从毛衣上的火花出发，一直走到“电场”这个看不见、摸不着，却能传递力和能量的概念。它是这本书里第一个真正意义上的“场”，也是后面所有电磁学、乃至理解光是什么的总入口。
\end{mainline}

\step{先猜一猜}

\begin{mainline}
往下读之前，先押三个注。请凭直觉回答，不必查资料，把答案写下来：

\begin{itemize}
  \item 用塑料尺在头发上摩擦几下，凑近桌上的碎纸屑，纸屑会被吸起来。那么纸屑被吸到尺子上以后，会一直老老实实粘在上面吗？
  \item 把两个气球都在头发上摩擦一遍，用线吊着靠近，它们会互相吸引、互相排斥，还是谁也不理谁？
  \item 摩擦让物体“带电”了。你觉得这些电是摩擦“制造”出来的，还是从某个地方“搬运”过来的？如果是制造出来的，制造之前它在哪里？
\end{itemize}

第三题看起来像是文字游戏，其实它直指本章最深的一条定律。至于第一题，几乎每个人都答错：大家默认“吸上去就粘住了”，而真实发生的现象比这个有趣得多——它里面藏着微观世界的一整套搬家过程。
\end{mainline}

\step{动手实验}

\begin{experiment}[一个气球的静电秀]
你需要：一个吹起来的气球，一些撕成指甲盖大小的碎纸屑，你的头发（或一件毛衣），一个能流出很细水流的水龙头。

第一步，吸纸屑。把气球在头发上用力摩擦十来下，然后悬停在纸屑上方两三厘米处。纸屑会跳起来，争先恐后地扑向气球。盯住看几秒钟：有些纸屑粘上气球之后，过一会儿又突然跳开、掉落。别急着翻解释，先把这个“先吸后弹”记在心里。

第二步，粘墙。把刚摩擦过的气球往墙上一按，松手。气球居然粘在墙上了，能挂几分钟。墙是平的、干的、不带任何胶水，也没有磁铁。是什么力顶住了气球的重力？

第三步，掰弯水流。把水龙头调到只流出一条细而连续的水线，把摩擦过的气球凑近水线（不要碰到）。水流会明显向气球一侧弯曲，像被一只看不见的手拉了一把。水既不是纸也不是铁，它为什么也会被“吸”？

第四步，互斥的两个气球。如果有两个气球，都摩擦一遍后分别用线吊起来靠近，你会看到它们明显地互相躲开。两个用同样方式处理过的气球，为什么偏偏要互相推离？

安全提示：这个实验完全安全，你感受到的电击能量极小。但请因此记住一个对应关系——加油站、面粉厂、喷漆车间里严禁静电火花，因为那里有可燃的油气与粉尘。同样是小火花，场合不同，后果天差地别。
\end{experiment}

\step{旧模型遇到麻烦}

\begin{mainline}
到目前为止，这本书里所有的力——推、拉、压、弹簧、摩擦、碰撞——都有一个共同点：需要接触。牛顿力学教我们画受力图，而受力图上的每一支箭头，都能找到一个“谁在碰它”的答案。就连万有引力虽然隔空，但你没法用它吸起纸屑：它太弱了，弱到只有地球这种庞然大物才刷得出存在感。现在拿这套观念来对付气球实验，会四处碰壁。

麻烦一：力隔空了。气球没有碰到纸屑，纸屑却跳了起来；气球没有碰到水流，水流却弯了。如果坚持“力必须有接触者”，我们就得发明一个看不见、却能推能拉的东西——这恰恰就是本章要干的事，但请先承认：旧模型里没给它留位置。

麻烦二：被吸的东西明明“没有电”。纸屑、水流、墙，谁都没有被摩擦过，按说“不带电”，可它们照样被吸引。如果吸引力是“带电物体之间”才有的，那不带电的东西凭什么被吸？更麻烦的是，两个都摩擦过的气球反而互相排斥。吸也吸得没道理，斥也斥得没道理。

麻烦三：电好像有两种。历史上人们发现，用丝绸摩擦玻璃棒得到的电，和用毛皮摩擦橡胶棒得到的电，行为不一样：两个玻璃棒互斥，两个橡胶棒互斥，玻璃棒和橡胶棒却互相吸引。旧模型里没有任何“两种”的余地——质量只有一种，引力只会吸引，从来不排斥。

麻烦四：账不平。仔细观察，摩擦从来不是单方面的：气球带电的同时，你的头发也带电，而且两者表现出来的“电性”恰好相反。电总是成对出现，像会计账本上的借和贷。如果电是被“制造”的，为什么每次不多不少恰好造出一对？

四条麻烦合起来，逼我们承认：需要一个新的量（它有两种、总量守恒），一种新的相互作用（隔空、可吸可斥、比引力强得惊人），以及一个新的传播者（那个看不见却能推手的东西）。
\end{mainline}

\step{发明新概念}

\begin{mainline}
像历史上的物理学家一样，我们一个一个地发明。

第一个概念：\textbf{电荷}。我们规定：物体带有一种属性，叫电荷；同种电荷互相排斥，异种电荷互相吸引。富兰克林给两种电荷起了名字：玻璃棒与丝绸摩擦后玻璃棒带的叫\textbf{正电荷}，橡胶棒与毛皮摩擦后橡胶棒带的叫\textbf{负电荷}。名字纯属约定——当初叫“甲电”“乙电”也行，但这个约定一旦定下来，全世界都要遵守。

第二个概念：\textbf{电荷守恒}。摩擦并不制造电荷，只是让电荷从一个物体搬到另一个物体。气球与头发摩擦，一些带负电的粒子从头发搬到了气球上：气球带负电，头发因失去它们而带等量的正电。一个孤立系统里，所有电荷的代数和保持不变——可以搬家，不能凭空生，也不能凭空灭。这回答了“先猜一猜”的第三题，也解释了麻烦四：电总是成对出现，因为每一次“带电”其实都是一次“分家”。这条定律至今没有任何例外，它和能量守恒、动量守恒一样，是物理学里最硬的规矩之一。后来人们发现，电荷的搬运工是\textbf{电子}——原子里带负电、最容易被抢走或塞进来的那一部分；原子核里的质子带正电，正常情况下被原子核牢牢扣住，不参与这种搬家。

值得停下来强调一句：上面这些不是空想出来的设定，而是被实验一条条逼出来的。最能说明问题的是密立根的油滴实验：让微小的油滴带上电，悬浮在两块金属板之间的电场里，调节电场让油滴恰好悬停不动，由平衡条件就能算出每颗油滴的电荷。结果干净利落——所有油滴的电荷都是同一个最小值的整数倍，从来没有出现过“半个台阶”。电荷像楼梯一样一级一级，这个最小的一级就是基本电荷 \(e\)，约 \(1.6\times10^{-19}\) 库仑。电荷量子化是测出来的事实，不是模型的假设。

第三个概念：\textbf{导体与绝缘体}。为什么金属门把手会电你，而毛衣不会主动电人？因为不同材料对电荷的态度截然不同。金属里有大量可以四处游荡的自由电子，电荷上去立刻散开、流动，这类材料叫\textbf{导体}；橡胶、塑料、干燥的头发里，电荷被锁在原地动弹不得，这类材料叫\textbf{绝缘体}。导体像铺平的马路，绝缘体像一格一格的抽屉。这个区别解释了大量日常现象：电缆要用铜芯外面裹塑料，一个负责导电、一个负责把电管住；冬天你的手是导体，一碰门把手，毛衣上攒的电荷瞬间经你的手泄放，于是“啪”地一下。

第四个概念用来解决麻烦二：\textbf{极化}。不带电的纸屑为什么会被吸引？纸屑内部正负电荷本来均匀混在一起，总电荷为零。带电气球靠近时，纸屑里的负电荷被排斥（或被吸引），在原子尺度上稍稍挪了挪位置——纸屑没有净电荷，却变成了“一头略正、一头略负”的小针。近的一头与气球异号，吸引力强；远的一头同号，排斥力弱；吸力与距离的依赖比斥力占优，净效果是吸引。水流被掰弯、气球粘墙，都是同一机制：\textbf{带电体不需要对方带电，只需要把对方的电荷“拨乱”一点}。这就是为什么带电物体能吸引轻小的中性物体——它不是例外，而是感应极化的必然结果。

这些概念还解决了一个历史悬案：为什么电总是出现在“摩擦”之后？其实关键不是摩擦这个动作，而是接触与分离。两种不同材料紧贴时，界面上对电子的“占有欲”不同，分开的瞬间一部分电子被留在了更“贪”的那一边。摩擦只是把接触面积放大、把接触次数增多，让转移的电子积少成多。干燥的冬天效果最明显，因为潮湿的空气会在物体表面铺一层薄薄的水膜，悄悄把积起来的电荷放走——所以静电是“干燥天气的特产”。

电荷守恒也不是一句空口宣言，它有精确的实验撑腰。法拉第做过一个著名的“冰桶实验”：把带电小球用丝线吊着，缓缓放进一个与静电计相连的金属桶里，不碰桶壁。小球一进桶，静电计立刻偏转；小球在桶里怎么晃动、怎么变换位置，读数纹丝不动；把小球提出来，读数回到零。桶外感应出的电荷，不多不少，恰好等于桶内小球携带的电荷——电荷既没有在感应过程中增殖，也没有被桶“吃掉”，它只是忠实地在内、外表面之间重新排队。后来更精密的实验把电荷守恒检验到了极高的精度，它至今屹立不倒。
\end{mainline}

\step{一句话模型}

\begin{oneline}
电荷是一种守恒的、有正负两种的属性，同性相斥、异性相吸；它通过充满空间的电场施加作用力——带电体先改变周围空间的“状态”，这个状态再去推或拉别的电荷，力不必接触就能传递。
\end{oneline}

\step{数学翻译器}

\begin{mainline}
概念齐了，现在把它们翻译成可以计算的式子。

第一条是力的大小。库仑用扭秤做了精密的测量：把带电小球挂在一根细得几乎不反抗扭转的石英丝下，让另一个带电小球靠近它，静电力把悬丝扭转一个角度，而扭转角度正比于力——用一把“会转身的天平”称出了比一粒灰尘还轻的力。结论干净利落：两个点电荷之间的力，与两个电荷量的乘积成正比，与距离的平方成反比，方向沿两者的连线——同号相斥、异号相吸。写成公式：

\[ F = k\,\frac{|q_1 q_2|}{r^2},\qquad k \approx 9.0\times 10^{9}\ \mathrm{N\cdot m^2/C^2} \]

这个式子回答“两个电荷之间的力有多大”。右边每一项都在改变什么：\(q_1 q_2\) 代表“双方各带多少电”，任一方加倍，力就加倍；\(r^2\) 在分母上，代表“离得越远，力按平方缩水”。电荷的国际单位叫库仑（C）。

公式一出，立刻用特殊情况检查它：

\begin{itemize}
  \item 距离拉大到无穷远，\(F\) 趋向零——隔得够远就互不相干，合理。
  \item 任一电荷为零，力为零——不带电就不参与（直接的）库仑作用，合理；纸屑那种情况靠的是极化“造出”局部电荷，不是违反这一条。
  \item 距离减半，力变为四倍——平方反比的脾气，这正是扭秤实验能精确检验的内容。
  \item 方向反转检查：交换两个电荷的位置，力大小不变、方向相反——正好是牛顿第三定律要求的“作用力与反作用力”，合理。
  \item 距离趋向零，力趋向无穷大——这是警报，不是预言：真实电荷不是数学上的点，凑近之后这个公式自己先失效，本章“模型的边界”一节会处理它。
\end{itemize}

不妨把它和万有引力放在一起看一眼。它像引力：都是平方反比，都是隔空作用的“长程力”。它不像引力：质量只有一种，引力永远吸引、无法屏蔽；电荷有两种，可以中和、可以屏蔽，而且强度悬殊到荒谬的程度——两个质子之间的电斥力，比它们之间的引力大约 \(10^{36}\) 倍。引力赢了宇宙，只是因为大尺度上正负电荷几乎完全抵消，引力却从不抵消。

第二条翻译是本章真正的主角：\textbf{电场}。库仑定律只说了“两个电荷互相推”，没说这个推是怎么隔着空气送到的。法拉第给出了一个改变物理学的想法：不要问 A 怎么推 B，要问 A 对它周围的空间做了什么。我们定义：空间每一点都有一个矢量，叫电场强度，大小等于“放在这一点的一库仑正电荷会受多大的力”，方向就是这个正电荷受力的方向：

\[ \mathbf E = \frac{\mathbf F}{q_{\text{试}}} \]

单位是牛每库仑（N/C）。这里 \(q_{\text{试}}\) 是一个想象中的\textbf{试探电荷}：足够小，小到它自己对原电场的扰动可以忽略。就像用一支不扰动水流的浮标去测水流。于是求力变成两步走：场源电荷产生电场，\(E = kQ/r^2\)（点电荷向外辐射或向内汇聚）；场再对场内电荷施力，\(\mathbf F = q\mathbf E\)。负电荷受力方向与场方向相反。

第三条是\textbf{叠加原理}：多个电荷共同产生的电场，等于每个电荷单独产生的电场按矢量相加。这条原理朴素得像废话，威力却巨大——它意味着再复杂的带电体，都可以拆成一小片一小片，分别算场再叠加。本章和第 20 章的几乎所有计算，骨子里都是这一句话。

电场看不见，我们给它画像：\textbf{电场线}。画一组曲线，线上每一点的切线方向就是该点电场的方向；线画得越密，表示场越强。点电荷的电场线是从电荷出发（或汇入电荷）的放射线，离得越远线越稀疏——而球面面积按 \(r^2\) 变大，线的密度恰好按 \(1/r^2\) 变稀，和库仑定律严丝合缝。这不是巧合，它暗示了一条更深的话，马上就说。
\end{mainline}

\begin{figure}[htbp]
\centering
\begin{tikzpicture}[scale=1]
  \fill[red!70!black] (0,0) circle (2.2pt) node[below=6pt] {$+Q$};
  \foreach \a in {0,45,...,315}
    \draw[->,thick] ({0.35*cos(\a)},{0.35*sin(\a)}) -- ({2.6*cos(\a)},{2.6*sin(\a)});
  \draw[dashed] (0,0) circle (1.1);
  \draw[dashed] (0,0) circle (2.2);
  \node at (3.5,1.9) {$r$ 加倍，};
  \node at (3.7,1.3) {线密度变 $1/4$};
\end{tikzpicture}
\caption{正点电荷的电场线：沿半径向外辐射，距离越远越稀疏，密度的下降正好是 \(1/r^2\)。}
\end{figure}

\begin{mainline}
现在做本章的估算，让你体会“库仑”是一个多么夸张的单位。设想两个小球各带一库仑电荷，相距一米：

\[ F = 9.0\times10^{9}\ \mathrm{N} \]

九十亿牛，约等于九十万吨物体的重量——差不多是两座摩天大楼压在一米的距离上。反过来算摩擦起电：一个气球经摩擦大约只转移了 \(10^{-8}\) 库仑量级、也就是约 \(10^{11}\) 个电子，而它表面的原子总数在 \(10^{23}\) 量级——平均每万亿个电子里只搬走一个，就足以让气球顶住重力粘在墙上、隔空掰弯水流。这两个估算合在一起，就是日常世界“电中性”的原因：电荷之间力的系数太大，任何宏观物体只要稍微积累一点净电荷，就会产生无法收拾的力，所以自然界拼命把正负电荷配平。静电现象之所以稀罕又醒目，恰恰因为它们是“配平工程”偶发的、极微量的失误。

最后介绍一条本章最漂亮、也最容易被当成“高级公式”跳过的规律——\textbf{高斯定律}的直觉版。想象电场是一种流体的流动，电荷是源头（或归宿）：正电荷像泉眼，不断向外“涌出”场线；负电荷像排水口，场线汇入消失。现在，不管你在泉眼外面套一个什么样的封闭面——球面、方盒子、土豆形状的怪面——\textbf{穿出这个面的场线“总量”，只由面里包着的净电荷决定，与面的形状、大小、位置全无关}。面包住的正电荷多，净流出就多；面里正负刚好抵消，净流出就是零。这个“穿出多少只认包住的电荷”的直觉，就是高斯定律。它像什么：像数一个房间里有多少人——只要在门口记录进出差额，不必追踪每个人的位置。它不像什么：它不像真的在数流体——电场不是任何物质的流动，“通量”只是对电场穿过曲面的一个定量度量，没有任何东西真的流过去。
\end{mainline}

\begin{mathbridge}
把上面的直觉写成式子，并看看它在对称情形下有多省力。定义电场通过封闭曲面的\textbf{电通量}为 \(\varPhi_E=\oint \mathbf E\cdot d\mathbf A\)（把曲面上每一小块的面积矢量与电场点乘，再全部加起来，出为正、入为负）。高斯定律说：

\[ \varPhi_E=\oint \mathbf E\cdot d\mathbf A=\frac{Q_{\text{内}}}{\varepsilon_0},\qquad \varepsilon_0\approx 8.85\times10^{-12}\ \mathrm{C^2/(N\cdot m^2)} \]

球对称：以点电荷 \(Q\) 为中心作半径 \(r\) 的球面。由对称性，球面上 \(E\) 处处大小相同且沿半径方向，于是 \(\varPhi_E = E\cdot 4\pi r^2 = Q/\varepsilon_0\)，解出

\[ E=\frac{Q}{4\pi\varepsilon_0 r^2} \]

——库仑定律的平方反比，原来是“通量守恒”加上“球面面积按 \(r^2\) 增长”的必然后果。换个角度说：我们住在三维空间里，这一条被写进了 \(1/r^2\) 里。

无限大均匀带电平面：由对称性，电场只能处处垂直于平面，且离平面多远都一样大。取一个穿过平面的圆柱面，两个底面各贡献 \(EA\)，侧面通量为零，故 \(2EA=\sigma A/\varepsilon_0\)，得

\[ E=\frac{\sigma}{2\varepsilon_0} \]

与距离无关——无限大平面的电场是均匀的。距离加倍场强不减半，这就是平方反比直觉遇到对称性时的反例，记得用特殊情况检查你的直觉。第 20 章的电容器，正是用两块这样的平面造出均匀电场的装置。
\end{mathbridge}

\begin{figure}[htbp]
\centering
\begin{tikzpicture}[scale=1]
  \draw[very thick] (-0.15,-1.6) -- (-0.15,1.6);
  \draw[very thick] (4.15,-1.6) -- (4.15,1.6);
  \foreach \y in {-1.2,-0.6,0,0.6,1.2}
    \node at (-0.38,\y) {$+$};
  \foreach \y in {-1.2,-0.6,0,0.6,1.2}
    \node at (4.38,\y) {$-$};
  \foreach \y in {-1.2,-0.6,0,0.6,1.2}
    \draw[->,thick] (0.25,\y) -- (3.75,\y);
  \node at (2,-2.1) {板间电场处处相同：同方向、同大小};
\end{tikzpicture}
\caption{两块带异号电荷的平行板之间，电场线平行等距——均匀电场。}
\end{figure}

\begin{mathbridge}

导体静电平衡也可以一句话拿下：导体内部若有电场，自由电子就会被推动，直到它们重新排布、把内部电场抵消为零。所以静电平衡时导体内部 \(E=0\)，电荷全住在表面上。金属外壳能屏蔽静电场，靠的就是这一条。
\end{mathbridge}

\begin{challenge}
高斯定律的积分形式还有一个局部的、微分的版本：\(\nabla\cdot\mathbf E=\rho/\varepsilon_0\)，读作“电场在某点的散度，等于该点的电荷密度除以 \(\varepsilon_0\)”。它把“泉眼在哪里”说成了一个逐点成立的等式，不再依赖任何想象中的封闭面。

再留一个为后文埋的伏笔。库仑定律写的是“此刻 A 在这里，B 受这样的力”，仿佛力的传递不需要时间。但如果 A 突然动了一下，B 会立刻知道吗？如果会，那就违反了一个更根本的限制——任何影响都不能超过光速传播（本书相对论部分会讲）。解决这个矛盾的办法，恰恰是把电场当成货真价实的东西：A 一动，扰动以光速在场中传开，B 要等扰动抵达才知道。一旦承认场自己会“携带消息”，第 22 章的磁场、第 24 章的麦克斯韦方程组、第 25 章的“光原来是电磁波”，就会一环扣一环地滚出来。电场是不是“真实存在”？到那一章你会发现，它携带能量和动量、可以被发射出去独立旅行——物理上对“真实”最苛刻的定义，它全都满足。
\end{challenge}

\step{模型的边界}

\begin{mainline}
老规矩，说清楚每个模型什么时候成立、什么时候失效，顺便分清“实验事实”“数学模型”和“物理解释”。

库仑定律的边界。它是对\textbf{静止}的点电荷（或球对称分布的电荷）的实验规律，在原子尺度到宏观尺度都被反复验证，精度极高。它的失效方式有三种。其一，电荷运动起来后，故事多出一半：运动的电荷还会感受到与速度有关的新力——磁力，那是第 22 章的内容，库仑定律只是完整规律的“静止部分”。其二，距离趋向零时公式发散，这暴露的是“点电荷”这个理想化：真实电荷总有一定分布，凑近了就得用叠加原理一小块一小块算。其三，电荷是量子化的：任何自由电荷都是基本电荷 \(e\approx1.6\times10^{-19}\) C 的整数倍。库仑定律里的 \(q\) 当成连续量用，在宏观上毫无问题，因为 \(e\) 太小；但讨论单个电子时，“半个电子的电荷”这种问法本身就错了。

电场线的边界。电场线是图，不是实体。空间里并没有一把一把的线；“线的密度表示场强”只是我们画图时的约定，而且只在特定画法下才严格对应通量。电场线有几条可以证明的脾气：不相交（否则交点处场有两个方向，矛盾）、从正电荷出发到负电荷终止或伸向无穷远、静电场中不闭合（这一点到第 23 章“变化的磁场制造电场”时会被打破，那时电场线可以闭合）。把电场线当成真实的弦，想象它们会“振动”“绷紧”，都是过度解读。

“场”这个解释的边界。实验事实是：电荷之间的力，大小方向如库仑定律所描述。数学模型是：库仑定律、叠加原理、高斯定律。物理解释是：“空间中存在电场，它传递相互作用”。三者别混淆：在你目前能做的所有静电实验里，光用“超距作用”也能算出同样的力，历史上也确实有人坚持超距观点。我们之所以认真地说“场是真实的东西”，是因为运动与变化的场合里超距观点会算错——变化的场携带能量、以有限速度传播，这些预言都被实验证实。这是“解释”凭新实验升格为“事实”的经典案例，也是本书区分三者的用意所在。

最后是两个常见的误用。第一，别把“试探电荷”当成一种新的电荷——它只是测量工具，像温度计不该加热被测的汤。第二，别把高斯定律当成“有对称性才能用的定律”——它对任何电荷分布、任何封闭面都严格成立；只是没有对称性时，你很难单靠它解出 \(E\)，需要叠加原理或更完整的方程。
\end{mainline}

\step{回头解释开场现象}

\begin{mainline}
现在回到那个冬夜的卧室，把小火花拆成完整的过程。

第一幕，分家。毛衣和头发是两种不同的材料，一整天里它们反复接触、摩擦，电子一点点从一方搬到另一方。毛衣上积累起净电荷——估算告诉我们，哪怕只有 \(10^{-8}\) 库仑量级，也已经足以制造可观的电场。冬天干燥，电荷没有水膜通道逃走，只能越攒越多。

第二幕，电场崛起。脱毛衣的动作让带电的两层突然分开。电荷已经分家，距离却在拉大，正负电荷之间的空间里电场随之增强。毛衣与头发之间狭小的间隙里，场强可以冲到每米几十万伏以上。

第三幕，击穿。空气平时是绝缘体，但这有个限度：场强达到大约 \(3\times10^{6}\) N/C（等价的说法是每米三百万伏）时，被电场拽着加速的零星自由电子会撞出更多电子，像雪崩一样，空气瞬间变成导体——这就是“击穿”。积累的电荷沿这条刚打通的通道泄放，能量以光和声的形式放出：你看到的蓝白火花、听到的“噼啪”，就是这个瞬间。

第四幕，举一反三。闪电是同一份剧本，只是演员换成了云：云内冰晶、水滴的对流与碰撞让电荷分层积累，云与地之间的电场增强，直到击穿几公里厚的空气，一次闪电转移的电荷可达十几库仑，释放的能量以十亿焦耳计。卧室火花与落地惊雷之间，差的只是规模，不是物理。

其余的开场现象也各就各位：气球粘墙，是气球上的净电荷把墙面极化，近端感应出异号电荷把它“扣”住；纸屑先吸后弹，是纸屑被极化后吸上去，接触气球时分得同号电荷，立刻从“异性相吸”变成“同性相斥”被弹开——第一题押“会一直粘着”的朋友，现在可以给自己打分了；门把手那一“啪”，是你身上积累的电荷经手指向金属泄放的微型击穿；水流弯曲，是水分子本身一头略正一头略负，在电场里集体转身、被拖向气球。

顺手还能解释一批你没问过、但天天见到的静电小事。老式显像管电视和电脑屏幕总是蒙一层灰，因为屏幕工作时带静电，把空气里的灰尘“钓”了过来；化纤裙子秋冬粘在腿上甩不开，是走路时布料与皮肤（或丝袜）反复摩擦、电荷分家后互相吸引；打印机偶尔会一次“吃进”好几张纸，往往就是干燥的纸张之间带了同号电荷粘叠在一起；理发店飘着的碎发会追着梳子走，和你开场的毛衣是同一回事。一旦眼里有了“电荷分家、极化、泄放”这三个动作，整个冬天都是你的实验室。

技术应用也都长在同一棵树上。静电除尘器让烟气先通过强电场使粉尘带电，再用带异号电荷的极板把它们从气流中“钓”出来，是大型工厂烟囱的标配；复印机和激光打印机用光在感光鼓上“写”出静电图案，墨粉只被吸到有电荷的地方，再转印到纸上——你手里的每张复印纸，都是一幅曾用电场作画的图。反向的工程同样重要：油罐车尾拖着铁链接地，加油站请你先摸一下释放静电，飞机降落前放掉机身电荷——都是在给“配平工程”帮忙，防止小火花出现在不该出现的地方。
\end{mainline}

\step{脑洞实验与挑战题}

\begin{thought}[如果质子和电子的电荷差一点点]
质子带正电、电子带负电，实验上两者电荷量相等到 \(10^{-21}\) 以上的精度，谁也没有测出过差别。现在开脑洞：假设电子电荷比质子小十亿分之一（\(10^{-9}\)），会发生什么？

你的身体里有约 \(10^{28}\) 量级的质子，配不平之后，每个质子平均多出 \(10^{-9}\) 个基本电荷的净正电，全身净电荷将达到约 \(1.6\) 库仑。回想本章的估算：相距一米的两个一库仑电荷，互斥九十亿牛。你和对面的人将各自带着库仑级的净电荷，互相以山崩般的力推离——别说握手，两个人根本无法站在同一个房间里。推而广之，地球也会带电，天体之间的电斥力将压过引力，星系、行星、原子层面的化学键全都改写。

这个脑洞要说的不是“好险”，而是一个深刻的事实：我们世界的存在本身，依赖于电荷极其精确的配平。为什么电子和质子的电荷绝对值恰好相等？这是物理学至今没有完全回答的问题之一，而大统一理论正是从这类问题里长出来的。
\end{thought}

\begin{puzzles}
  \item 带电气球吸起纸屑后，纸屑过几秒会跳开。有人说“是因为气球的电用完了”。这个说法错在哪里？请用电荷守恒和接触带电重新讲一遍全过程。\\
  \textbf{思路提示：}电荷不会“用完”，只会搬家。追踪电子从气球到纸屑的那一步：纸屑吸上去之前带什么电？接触之后带什么电？
  \item 两个相同的带电气球固定在桌面对称位置，第三个带电小球用线吊在它们正中间。第三个小球受的合力为零。现在把它轻轻往旁边推一小段再松手，它会回到中间，还是越走越远？这个“合力为零的位置”和碗底的小球有什么不同？\\
  \textbf{思路提示：}画出它偏离中心后两个库仑力的方向与大小变化——离谁近，谁的力按平方变大。比较“把它拉回去的力”与“把它推得更远的力”谁赢。这其实是一条普遍定理（静电场中不存在仅靠静电力稳定的平衡点）的直觉版本。
  \item 把正在播放的收音机用金属网罩整个罩住，声音立刻变弱甚至消失；用塑料网罩住则几乎没有影响。用本章“导体静电平衡”的知识解释金属网罩做了什么，并说说为什么网眼不必做得像铁皮那样严丝合缝，也能起相当作用。\\
  \textbf{思路提示：}导体内部的电场会被自由电荷的重新排布抵消。再想一想：罩内场被削弱的前提，是网孔尺寸远小于想要挡住的那个变化的“空间尺度”——这条线索会在第 25 章变成严格的结论。
  \item 本章说高斯定律的直觉是“穿出封闭面的总量只认面内的电荷”。假如库仑定律不是平方反比，而是按 \(1/r^3\) 衰减，这个直觉还成立吗？用点电荷外套两个不同半径的球面检验一下。\\
  \textbf{思路提示：}通量等于场强乘以球面面积。若 \(E\) 按 \(1/r^3\) 衰减，内球面与外球面的通量各是多少？它们还相等吗？由此你会看到：平方反比不是库仑定律的巧合，而是“通量守恒”能成立的前提——或者说，是三维空间的签名。
\end{puzzles}
